\documentclass[12pt, titlepage]{article}

\usepackage{fullpage}
\usepackage[round]{natbib}
\usepackage{multirow}
\usepackage{booktabs}
\usepackage{tabularx}
\usepackage{graphicx}
\usepackage{float}
\usepackage{hyperref}
\hypersetup{
    colorlinks,
    citecolor=blue,
    filecolor=black,
    linkcolor=red,
    urlcolor=blue
}

\input{../../Comments.text}
\input{../../Common.text}

\newcounter{acnum}
\newcommand{\actheacnum}{AC\theacnum}
\newcommand{\acref}[1]{AC\ref{#1}}

\newcounter{ucnum}
\newcommand{\uctheucnum}{UC\theucnum}
\newcommand{\uref}[1]{UC\ref{#1}}

\newcounter{mnum}
\newcommand{\mthemnum}{M\themnum}
% Module numbers. \mitem{key} numbers a module in the hierarchy table and
% links the number to the module's entry, labelled sec:key, in the Module
% Decomposition section. \mref{key} prints the number with the same link.
\newcommand{\mref}[1]{\hyperref[sec:#1]{M\ref*{#1}}}
\newcommand{\mitem}[1]{\refstepcounter{mnum}\label{#1}\mref{#1}}

\begin{document}

\title{Module Guide for \progname{}} 
\author{\authname}
\date{\today}

\maketitle

\pagenumbering{roman}

\section{Revision History}

\begin{tabularx}{\textwidth}{p{3cm}p{2cm}X}
\toprule {\bf Date} & {\bf Version} & {\bf Notes}\\
\midrule
Date 1 & 1.0 & Notes\\
Date 2 & 1.1 & Notes\\
\bottomrule
\end{tabularx}

\newpage

\section{Reference Material}

This section records information for easy reference. 

\subsection{Abbreviations and Acronyms}

\renewcommand{\arraystretch}{1.2}
\begin{tabular}{l l} 
  \toprule		
  \textbf{symbol} & \textbf{description}\\
  \midrule 
  AC & Anticipated Change\\
  DAG & Directed Acyclic Graph \\
  M & Module \\
  MG & Module Guide \\
  OS & Operating System \\
  R & Requirement\\
  SC & Scientific Computing \\
  SRS & Software Requirements Specification\\
  \progname & Distributed system for local artificial intelligence inference. \\
  UC & Unlikely Change \\
  \wss{etc.} & \wss{...}\\
  \bottomrule
\end{tabular}\\

\newpage

\tableofcontents

\listoftables

\listoffigures

\newpage

\pagenumbering{arabic}

\section{Introduction}

Decomposing a system into modules is a commonly accepted approach to developing
software.  A module is a work assignment for a programmer or programming
team~\citep{ParnasEtAl1984}.  We advocate a decomposition
based on the principle of information hiding~\citep{Parnas1972a}.  This
principle supports design for change, because the ``secrets'' that each module
hides represent likely future changes.  Design for change is valuable in SC,
where modifications are frequent, especially during initial development as the
solution space is explored.  

Our design follows the rules layed out by \citet{ParnasEtAl1984}, as follows:
\begin{itemize}
\item System details that are likely to change independently should be the
  secrets of separate modules.
\item Each data structure is implemented in only one module.
\item Any other program that requires information stored in a module's data
  structures must obtain it by calling access programs belonging to that module.
\end{itemize}

After completing the first stage of the design, the Software Requirements
Specification (SRS), the Module Guide (MG) is developed~\citep{ParnasEtAl1984}. The MG
specifies the modular structure of the system and is intended to allow both
designers and maintainers to easily identify the parts of the software.  The
potential readers of this document are as follows:

\begin{itemize}
\item New project members: This document can be a guide for a new project member
  to easily understand the overall structure and quickly find the
  relevant modules they are searching for.
\item Maintainers: The hierarchical structure of the module guide improves the
  maintainers' understanding when they need to make changes to the system. It is
  important for a maintainer to update the relevant sections of the document
  after changes have been made.
\item Designers: Once the module guide has been written, it can be used to
  check for consistency, feasibility, and flexibility. Designers can verify the
  system in various ways, such as consistency among modules, feasibility of the
  decomposition, and flexibility of the design.
\end{itemize}

The rest of the document is organized as follows. Section
\ref{SecChange} lists the anticipated and unlikely changes of the software
requirements. Section \ref{SecMH} summarizes the module decomposition that
was constructed according to the likely changes. Section \ref{SecConnection}
specifies the connections between the software requirements and the
modules. Section \ref{SecMD} gives a detailed description of the
modules. Section \ref{SecTM} includes two traceability matrices. One checks
the completeness of the design against the requirements provided in the SRS. The
other shows the relation between anticipated changes and the modules. Section
\ref{SecUse} describes the use relation between modules.

\section{Anticipated and Unlikely Changes} \label{SecChange}

This section lists possible changes to the system. According to the likeliness
of the change, the possible changes are classified into two
categories. Anticipated changes are listed in Section \ref{SecAchange}, and
unlikely changes are listed in Section \ref{SecUchange}.

\subsection{Anticipated Changes} \label{SecAchange}

Anticipated changes are the source of the information that is to be hidden
inside the modules. Ideally, changing one of the anticipated changes will only
require changing the one module that hides the associated decision. The approach
adapted here is called design for
change.

\begin{description}
\item[\refstepcounter{acnum} \actheacnum \label{acHardware}:] The specific
  hardware on which the software is running.
\item[\refstepcounter{acnum} \actheacnum \label{acInput}:] The format of the
  initial input data.
\item[\refstepcounter{acnum} \actheacnum \label{acUI}:] The layout and visual design 
  of the Lake Client user interface (network graph, metrics
  display, and chat view). These will likely change based on user feedback.
\item[\refstepcounter{acnum} \actheacnum \label{acMetrics}:] The set of device and
  network health and performance metrics displayed to the user, and how they are
  presented (e.g., bottleneck highlighting).
\item[\refstepcounter{acnum} \actheacnum \label{acCoordAPI}:] The request and
  response format of the Coordinator's REST API, as consumed by the user
  interface.
\end{description}

\subsection{Unlikely Changes} \label{SecUchange}

The module design should be as general as possible. However, a general system is
more complex. Sometimes this complexity is not necessary. Fixing some design
decisions at the system architecture stage can simplify the software design. If
these decision should later need to be changed, then many parts of the design
will potentially need to be modified. Hence, it is not intended that these
decisions will be changed.

\begin{description}
\item[\refstepcounter{ucnum} \uctheucnum \label{ucIO}:] Input/Output devices accepted
by the system will not change from the following: Input - File, Keyboard, Mouse, and Output - 
File, Keyboard, Mouse.
\item[\refstepcounter{ucnum} \uctheucnum \label{ucREST}:] The REST API host that the user
interface will rely on is unlikely to change, since the Coordinator module is responsible for handling all external APIs. 
\item[\refstepcounter{ucnum} \uctheucnum \label{ucOtherAPI}:] The Lake Client will only contact the Coordinator for
information and requests in the system. Any additional APIs beyond the one between the interface and the Coordinator 
are unlikely to be added.
\item[\refstepcounter{ucnum} \uctheucnum \label{ucFEFramework}:] The Lake Client user interface's frontend framework is highly unlikely to change.
This is due to the significant rewriting effort that would be required to change the module from one framework to another.
\end{description}

\section{Module Hierarchy} \label{SecMH}

This section provides an overview of the module design. Modules are summarized
in a hierarchy decomposed by secrets in Table \ref{TblMH}. The modules listed
below, which are leaves in the hierarchy tree, are the modules that will
actually be implemented.

\begin{description}
\item [\mitem{mGGML}:] GGML Module
\item [\mitem{mCoordinator}:] Coordinator Module
\item [\mitem{mCore}:] Core Module
\item [\mitem{mIE}:] Inference Engine Module
\item [\mitem{mLC}:] Lake Client (User Interface) Module
\item [\mitem{mPartition}:] Partition Module
\item [\mitem{mTIO}:] Tensor I/O Module
\end{description}

\begin{table}[h!]
\centering
\begin{tabular}{p{0.3\textwidth} p{0.6\textwidth}}
\toprule
\textbf{Level 1} & \textbf{Level 2}\\
\midrule

{Hardware-Hiding Module} & \mref{mGGML}: GGML Module\\
\midrule

\multirow{4}{0.3\textwidth}{Behaviour-Hiding Module} & \mref{mCoordinator}: Coordinator Module\\
& \mref{mCore}: Core Module\\
& \mref{mIE}: Inference Engine Module\\
& \mref{mLC}: Lake Client Module\\
\midrule

\multirow{2}{0.3\textwidth}{Software Decision Module} & \mref{mPartition}: Partition Module\\
& \mref{mTIO}: Tensor I/O Module\\
\bottomrule

\end{tabular}
\caption{Module Hierarchy}
\label{TblMH}
\end{table}

\section{Connection Between Requirements and Design} \label{SecConnection}

The design of the system is intended to satisfy the requirements developed in
the SRS. In this stage, the system is decomposed into modules. The connection
between requirements and modules is listed in Table~\ref{TblRT}.

The following design decisions connect the Lake Client module (\mref{mLC}) to the SRS:
\begin{itemize}
\item \textbf{Device discovery, selection, and model choice} (FR.S.5.1, FR.S.5.2, FR.S.5.3): The Lake Client lists the devices found by the Coordinator, lets the user choose which to include, and lets the user pick a GGUF model file. It passes these choices to the Coordinator as network requests and does not communicate with other nodes itself.
\item \textbf{Infeasible partitioning} (FR.S.5.4): The Model Partitioner reports a failure through the Coordinator, and the Lake Client displays the error to the user.
\item \textbf{Metrics and bottlenecks} (FR.S.5.5, FR.S.5.6): The Lake Client only renders the latest metrics snapshot provided by the Coordinator, for the network, each node, and each connection, and highlights the bottleneck the Coordinator identifies. It does not store metric history.
\item \textbf{Usability} (NFR.S.1, NFR.S.2, NFR.S.7): The layout and visual design are hidden inside this module (\acref{acUI}) so they can be revised after usability testing without affecting other modules.
\item \textbf{Chat interface}: The chat view is a convenience for starting and continuing inference sessions from the Lake Client. All inference requests are routed through the Coordinator, which serves the OpenAI-compatible API (FR.S.1.9).
\end{itemize}

\section{Module Decomposition} \label{SecMD}

Modules are decomposed according to the principle of ``information hiding''
proposed by \citet{ParnasEtAl1984}. The \emph{Secrets} field in a module
decomposition is a brief statement of the design decision hidden by the
module. The \emph{Services} field specifies \emph{what} the module will do
without documenting \emph{how} to do it. For each module, a suggestion for the
implementing software is given under the \emph{Implemented By} title. If the
entry is \emph{OS}, this means that the module is provided by the operating
system or by standard programming language libraries.  \emph{\progname{}} means the
module will be implemented by the \progname{} software.

Only the leaf modules in the hierarchy have to be implemented. If a dash
(\emph{--}) is shown, this means that the module is not a leaf and will not have
to be implemented.

\subsection{Hardware Hiding Modules (\mref{mHH})}

\begin{description}
\item[Secrets:]The data structure and algorithm used to implement the virtual
  hardware.
\item[Services:]Serves as a virtual hardware used by the rest of the
  system. This module provides the interface between the hardware and the
  software. So, the system can use it to display outputs or to accept inputs.
\item[Implemented By:] OS
\end{description}

\subsection{Behaviour-Hiding Module}

\begin{description}
\item[Secrets:]The contents of the required behaviours.
\item[Services:]Includes programs that provide externally visible behaviour of
  the system as specified in the software requirements specification (SRS)
  documents. This module serves as a communication layer between the
  hardware-hiding module and the software decision module. The programs in this
  module will need to change if there are changes in the SRS.
\item[Implemented By:] --
\end{description}

\subsubsection{Coordinator Module (\mref{mCoordinator})} \label{sec:mCoordinator}

\begin{description}
\item[Secrets:] How the Lake network is formed, monitored and rebuilt, and how
  inference requests enter and leave it.
\end{description}

\subsubsection{Core Module (\mref{mCore})} \label{sec:mCore}

\begin{description}
\item[Secrets:] How a node joins a Lake network, holds its partition and
  reports its health.
\end{description}

\subsubsection{Lake Client (User Interface) (\mref{mLC})} \label{sec:mLC}

\begin{description}
\item[Secrets:] Implementation of design choices with regards to user interfaces and user experience,
routing of requests and underlying services.
\end{description}

\subsection{Software Decision Module}

\begin{description}
\item[Secrets:] The design decision based on mathematical theorems, physical
  facts, or programming considerations. The secrets of this module are
  \emph{not} described in the SRS.
\item[Services:] Includes data structure and algorithms used in the system that
  do not provide direct interaction with the user. 
  % Changes in these modules are more likely to be motivated by a desire to
  % improve performance than by externally imposed changes.
\item[Implemented By:] --
\end{description}

\subsubsection{Etc.}

\section{Traceability Matrix} \label{SecTM}

This section shows two traceability matrices: between the modules and the
requirements and between the modules and the anticipated changes.

% the table should use mref, the requirements should be named, use something
% like fref
\begin{table}[H]
\centering
\begin{tabular}{p{0.2\textwidth} p{0.6\textwidth}}
\toprule
\textbf{Req.} & \textbf{Modules}\\
\midrule
FR.S.5.1, FR.S.5.2, FR.S.5.3 & \mref{mUserInterface}\\
FR.S.5.4 & \mref{mUserInterface}\\
FR.S.5.5, FR.S.5.6 & \mref{mUserInterface}\\
NFR.S.1, NFR.S.2, NFR.S.7 & \mref{mUserInterface}\\
% Add rows for the other modules here.
\bottomrule
\end{tabular}
\caption{Trace Between Requirements and Modules}
\label{TblRT}
\end{table}

\begin{table}[H]
\centering
\begin{tabular}{p{0.2\textwidth} p{0.6\textwidth}}
\toprule
\textbf{AC} & \textbf{Modules}\\
\midrule
\acref{acHardware} & \mref{mHH}\\
\acref{acInput} & \mref{mInput}\\
\acref{acParams} & \mref{mParams}\\
\acref{acVerify} & \mref{mVerify}\\
\acref{acOutput} & \mref{mOutput}\\
\acref{acVerifyOut} & \mref{mVerifyOut}\\
\acref{acODEs} & \mref{mODEs}\\
\acref{acEnergy} & \mref{mEnergy}\\
\acref{acControl} & \mref{mControl}\\
\acref{acSeqDS} & \mref{mSeqDS}\\
\acref{acSolver} & \mref{mSolver}\\
\acref{acPlot} & \mref{mPlot}\\
\bottomrule
\end{tabular}
\caption{Trace Between Anticipated Changes and Modules}
\label{TblACT}
\end{table}

\section{Use Hierarchy Between Modules} \label{SecUse}

In this section, the uses hierarchy between modules is
provided. \citet{Parnas1978} said of two programs A and B that A {\em uses} B if
correct execution of B may be necessary for A to complete the task described in
its specification. That is, A {\em uses} B if there exist situations in which
the correct functioning of A depends upon the availability of a correct
implementation of B.  Figure \ref{FigUH} illustrates the use relation between
the modules. It can be seen that the graph is a directed acyclic graph
(DAG). Each level of the hierarchy offers a testable and usable subset of the
system, and modules in the higher level of the hierarchy are essentially simpler
because they use modules from the lower levels.

\wss{The uses relation is not a data flow diagram.  In the code there will often
be an import statement in module A when it directly uses module B.  Module B
provides the services that module A needs.  The code for module A needs to be
able to see these services (hence the import statement).  Since the uses
relation is transitive, there is a use relation without an import, but the
arrows in the diagram typically correspond to the presence of import statement.}

\wss{If module A uses module B, the arrow is directed from A to B.}

\begin{figure}[H]
\centering
%\includegraphics[width=0.7\textwidth]{UsesHierarchy.png}
\caption{Use hierarchy among modules}
\label{FigUH}
\end{figure}

%\section*{References}

\section{User Interfaces}

\wss{Design of user interface for software and hardware.  Attach an appendix if
needed. Drawings, Sketches, Figma}

\section{Design of Communication Protocols}

\wss{If there is communication between hardware and software, it may be necessary 
to explain the design of the communication protocol.  This section should be filled 
in with NA if it is not appropriate for a given project.}

\section{Database Design}

\wss{If a database is a significant component of the project, the database
design should be summarized here.  This section should be filled in with NA if
it is not appropriate for a given project.}

\section{Timeline}

\wss{Schedule of tasks and who is responsible}

\wss{You can point to GitHub if this information is included there}

\section{Generative AI Use Disclosure}

\subsection{AI Tools Used}

Claude Sonnet 5.5

\subsection{How and Where Used}

Sonnet 5.5 was used to assist in the mapping between the Lake Client (User Interface) module and the SRS contents for the connection between requirements and design, along with the traceability table.

\subsection{Verification of AI-Generated Content}

Any generated content was proofread and verified to confirm accuracy and alignment with existing team objectives and documentation.

\newpage{}

\bibliographystyle {plainnat}
\bibliography{../../../refs/References}

\end{document}